\documentclass[a4paper,12pt]{article}
\usepackage[utf8]{inputenc}
\usepackage{amsmath, amssymb, amsthm}
\usepackage{geometry}
\geometry{top=2cm, bottom=2cm, left=2cm, right=2cm}
\usepackage[T2A]{fontenc}
\usepackage[english,russian]{babel}

\newtheorem{definition}{Определение}
\newtheorem{theorem}{Теорема}
\newtheorem{corollary}{Следствие}
\theoremstyle{definition}
\newtheorem{example}{Пример}
\renewcommand{\proofname}{Доказательство}

\begin{document}

\title{Подробный конспект ответов на экзаменационные вопросы\\
\large{по курсу «Обыкновенные дифференциальные уравнения»}}
\author{По учебнику В.И. Арнольда}
\date{}
\maketitle

\section{Основные понятия}

\begin{definition}[Гладкое отображение]
	Отображение $f:U\subset\mathbb{R}^m\to\mathbb{R}^n$ называется \emph{гладким} (класса $C^r$), если все его частные производные до порядка $r$ включительно существуют и непрерывны. При $r=\infty$ говорят о бесконечно гладком отображении.
\end{definition}

\begin{definition}[Гладкое векторное поле]
	\emph{Векторным полем} в области $U\subset\mathbb{R}^n$ называется отображение $v:U\to\mathbb{R}^n$, сопоставляющее каждой точке $x\in U$ вектор $v(x)$.
	Поле называется \emph{гладким}, если оно является гладким отображением.
	Векторы $v(x)$ интерпретируются как скорости движения фазовой точки.
\end{definition}

\begin{definition}[Обыкновенное дифференциальное уравнение первого порядка]
	Пусть $v(t,x)$ – гладкое векторное поле, заданное в области расширенного фазового пространства $\mathbb{R}\times\mathbb{R}^n$.
	Уравнение
	\[
	\dot{x}=v(t,x),\qquad x\in\mathbb{R}^n,\; t\in\mathbb{R},
	\]
	называется \emph{обыкновенным дифференциальным уравнением первого порядка}.
	Если $v$ не зависит от $t$, уравнение называется \emph{автономным} (однородным по времени).
\end{definition}

\begin{definition}[Фазовое пространство]
	Множество всех возможных состояний системы называется \emph{фазовым пространством} (обычно $\mathbb{R}^n$).
	Прямое произведение оси времени на фазовое пространство $\mathbb{R}\times\mathbb{R}^n$ называют \emph{расширенным фазовым пространством}.
\end{definition}

\begin{definition}[Решение, интегральная кривая]
	Гладкое отображение $\varphi:I\to\mathbb{R}^n$ ($I$ – интервал) называется \emph{решением} уравнения $\dot{x}=v(t,x)$, если
	\[
	\frac{d\varphi}{dt}=v(t,\varphi(t))\quad\forall t\in I.
	\]
	График решения в расширенном фазовом пространстве называется \emph{интегральной кривой}.
\end{definition}

\section{Однородные и неоднородные уравнения первого порядка}

\subsection{Линейные уравнения}

\begin{definition}[Линейное однородное уравнение]
	Уравнение
	\[
	\frac{dx}{dt}=a(t)x,\qquad x\in\mathbb{R},
	\]
	называется \emph{линейным однородным уравнением первого порядка}.
	Его общее решение: $x(t)=C\exp\!\Bigl(\int a(t)\,dt\Bigr)$.
\end{definition}

\begin{example}[Экспоненциальный рост]
	Уравнение $\dot{x}=kx$ (нормальное размножение). Решение: $x(t)=x_0e^{kt}$.
	При $k>0$ – рост, при $k<0$ – убывание.
\end{example}

\begin{definition}[Линейное неоднородное уравнение]
	Уравнение
	\[
	\frac{dx}{dt}=a(t)x+b(t)
	\]
	называется \emph{линейным неоднородным уравнением первого порядка}.
\end{definition}

\begin{theorem}[Формула для решения неоднородного уравнения]
	Пусть $x_0(t)=\exp\!\bigl(\int_{t_0}^t a(\tau)d\tau\bigr)$ – ненулевое решение однородного уравнения.
	Тогда решение неоднородного уравнения с начальным условием $x(t_0)=x_0$ даётся формулой
	\[
	x(t)=x_0(t)\Bigl(x_0+\int_{t_0}^t x_0^{-1}(\tau)b(\tau)\,d\tau\Bigr).
	\]
\end{theorem}
\begin{proof}[Метод вариации постоянной]
	Ищем решение в виде $x(t)=C(t)x_0(t)$. Подстановка в уравнение даёт $\dot C x_0+C\dot x_0 = a C x_0 + b$. Учитывая, что $\dot x_0 = a x_0$, получаем $\dot C x_0 = b$, откуда $\dot C = b/x_0$. Интегрируя: $C(t)=C_0+\int_{t_0}^t b/x_0\,d\tau$. Подставляя обратно, имеем $x(t)=x_0(t)\bigl(C_0+\int_{t_0}^t b/x_0\,d\tau\bigr)$. Начальное условие даёт $C_0=x_0$.
\end{proof}

\subsection{Однородные (в смысле однородности функции) уравнения}

\begin{definition}[Однородное уравнение]
	Уравнение
	\[
	\frac{dy}{dx}=F\!\left(\frac{y}{x}\right)
	\]
	называется \emph{однородным}. Оно инвариантно относительно растяжений $(x,y)\mapsto(e^s x,e^s y)$.
\end{definition}

\begin{theorem}[Интегрирование однородного уравнения]
	Замена $y=vx$ (при $x>0$) приводит однородное уравнение к уравнению с разделяющимися переменными:
	\[
	x\frac{dv}{dx}=F(v)-v.
	\]
\end{theorem}

\begin{example}
	Решим уравнение $\dfrac{dy}{dx}=\dfrac{y}{x}+\dfrac{y^2}{x^2}$.
	Положим $y=vx$. Тогда $\dfrac{dy}{dx}=v+x\dfrac{dv}{dx}$. Подставляя: $v+x\dfrac{dv}{dx}=v+v^2$, откуда $x\dfrac{dv}{dx}=v^2$.
	Разделяем переменные: $\dfrac{dv}{v^2}=\dfrac{dx}{x}$. Интегрируем: $-\dfrac{1}{v}=\ln x+C$, т.е. $v=-\dfrac{1}{\ln x+C}$. Возвращаясь к $y$: $y=-\dfrac{x}{\ln x+C}$.
\end{example}

\subsection{Квазиоднородные уравнения}

\begin{definition}[Квазиоднородное уравнение]
	Уравнение $\dfrac{dy}{dx}=F(x,y)$ называется \emph{квазиоднородным} с весами $\alpha=\deg x$, $\beta=\deg y$, если его поле направлений инвариантно относительно растяжений $(x,y)\mapsto(e^{\alpha s}x,e^{\beta s}y)$.
	Для этого необходимо, чтобы $F$ была квазиоднородной степени $\beta-\alpha$.
\end{definition}

\begin{theorem}[Интегрирование квазиоднородного уравнения]
	Подстановка $v=y^\alpha/x^\beta$ (в области $x>0$) сводит квазиоднородное уравнение к уравнению с разделяющимися переменными.
\end{theorem}

\begin{example}
	Уравнение $\dfrac{dy}{dx}=-\dfrac{x}{y^3}$.
	Подберём веса: пусть $\deg x=2$, $\deg y=1$. Тогда $\deg y^3=3$, $\deg x=2$, и правая часть имеет степень $2-3=-1$, а левая $\dfrac{dy}{dx}$ имеет степень $1-2=-1$. Условие выполнено.
	Положим $v=y/x^{1/2}$ (т.к. $\beta/\alpha=1/2$). Тогда $y=v\sqrt{x}$, $dy=\sqrt{x}\,dv+\frac{v}{2\sqrt{x}}dx$. Подставляя и упрощая, получим уравнение с разделяющимися переменными.
\end{example}

\section{Устойчивость решений ОДУ}

\subsection{Устойчивость по Ляпунову и асимптотическая устойчивость}

Рассмотрим автономное уравнение
\[
\dot{x}=v(x),\qquad x\in\mathbb{R}^n,
\]
с положением равновесия $x=0$ ($v(0)=0$).

\begin{definition}[Устойчивость по Ляпунову]
	Положение равновесия $x=0$ называется \emph{устойчивым по Ляпунову}, если для любого $\varepsilon>0$ найдётся $\delta>0$ такое, что для любого начального условия $x_0$ с $|x_0|<\delta$ решение $\varphi(t)$ с $\varphi(0)=x_0$ определено при всех $t\ge0$ и $|\varphi(t)|<\varepsilon$ для всех $t\ge0$.
\end{definition}

\begin{definition}[Асимптотическая устойчивость]
	Положение равновесия $x=0$ называется \emph{асимптотически устойчивым}, если оно устойчиво по Ляпунову и, кроме того,
	\[
	\lim_{t\to+\infty}\varphi(t)=0
	\]
	для всех решений с начальными условиями, достаточно близкими к $0$.
\end{definition}

\subsection{Функция Ляпунова и критерий устойчивости}

\begin{definition}[Функция Ляпунова]
	Пусть $U$ – окрестность точки $0$, $V:U\to\mathbb{R}$ – гладкая функция, $V(0)=0$, $V(x)>0$ при $x\neq0$ (положительно определена).
	Если производная $V$ вдоль поля $v$ (т.е. $L_v V = \sum\frac{\partial V}{\partial x_i}v_i$) удовлетворяет условию
	\[
	L_v V(x) \le 0 \quad (\text{или } L_v V(x) < 0 \text{ при } x\neq0),
	\]
	то $V$ называется \emph{функцией Ляпунова} для системы $\dot{x}=v(x)$.
\end{definition}

\begin{theorem}[Критерий устойчивости (Ляпунов)]
	\begin{enumerate}
		\item Если в некоторой окрестности положения равновесия $x=0$ существует положительно определённая функция $V$ такая, что $L_v V\le 0$, то положение равновесия устойчиво по Ляпунову.
		\item Если же $L_v V <0$ для всех $x\neq0$ (в окрестности), то равновесие асимптотически устойчиво.
	\end{enumerate}
\end{theorem}
\begin{proof}[Доказательство]
	1) Фиксируем $\varepsilon>0$ настолько малым, что шар $|x|\le\varepsilon$ лежит в $U$. Положим $m=\min_{|x|=\varepsilon} V(x)>0$ (в силу положительной определённости). Выберем $\delta>0$ так, чтобы при $|x|<\delta$ выполнялось $V(x)<m$ (это возможно, так как $V(0)=0$ и $V$ непрерывна). Возьмём любое решение $\varphi(t)$ с $|\varphi(0)|<\delta$. Поскольку $L_v V\le 0$, функция $V(\varphi(t))$ не возрастает. Следовательно, $V(\varphi(t))\le V(\varphi(0))<m$ для всех $t\ge0$. Если бы в какой-то момент $t_1$ оказалось $|\varphi(t_1)|=\varepsilon$, то $V(\varphi(t_1))\ge m$, противоречие. Значит, $|\varphi(t)|<\varepsilon$ для всех $t\ge0$. Устойчивость доказана.
	
	2) Пусть $L_v V<0$ при $x\neq0$. Тогда $V$ строго убывает вдоль траекторий. Покажем, что $V(\varphi(t))\to0$. Предположим противное: $V(\varphi(t))\ge c>0$ для всех $t$. Тогда в силу строгого убывания существует предел $c_0>0$. Рассмотрим компактное множество $V(x)\ge c_0$, $|x|\le\varepsilon$. На нём $L_v V\le -\alpha<0$ (непрерывная функция достигает максимума). Тогда $\frac{d}{dt}V(\varphi(t))\le -\alpha$, что приводит к $V(\varphi(t))\le V(\varphi(0))-\alpha t\to -\infty$ – противоречие. Следовательно, $V(\varphi(t))\to0$. Так как $V$ положительно определена, $\varphi(t)\to0$.
\end{proof}

\subsection{Теорема об устойчивости по первому приближению}

Пусть $v(0)=0$. Линеаризованная система: $\dot{x}=Ax$, где $A=v_*(0)$ – матрица Якоби.

\begin{theorem}[Об устойчивости по первому приближению]
	Если все собственные числа $A$ имеют отрицательные вещественные части, то положение равновесия $x=0$ асимптотически устойчиво.
\end{theorem}
\begin{proof}[Доказательство]
	Из теории линейных систем (см. §22) известно, что при $\operatorname{Re}\lambda<0$ существует квадратичная форма $V(x)=|x|^2$ в подходящих координатах, для которой $L_{Ax}V$ – отрицательно определённая форма, т.е. $L_{Ax}V\le -\gamma|x|^2$, $\gamma>0$.
	Исходное поле $v(x)=Ax+o(|x|)$. Тогда
	\[
	L_v V = L_{Ax}V + L_{o(|x|)}V.
	\]
	Оценим второе слагаемое: $L_{o(|x|)}V = \sum 2x_i\cdot o(|x|)=O(|x|^3)$. Следовательно, при достаточно малых $|x|$ имеем
	\[
	L_v V \le -\gamma|x|^2 + C|x|^3 \le -\frac{\gamma}{2}|x|^2,
	\]
	если $|x|$ достаточно мал (например, $|x|<\gamma/(2C)$). Таким образом, $V$ является функцией Ляпунова с $L_v V<0$ в некоторой окрестности $0$. По критерию Ляпунова положение равновесия асимптотически устойчиво.
\end{proof}

\begin{example}[Маятник с трением]
	Уравнение малых колебаний маятника с трением: $\ddot{x}+k\dot{x}+x=0$ ($k>0$). Система первого порядка: $\dot{x}_1=x_2$, $\dot{x}_2=-x_1-kx_2$. Матрица $A=\begin{pmatrix}0&1\\-1&-k\end{pmatrix}$ имеет собственные числа $\lambda_{1,2}=-\frac{k}{2}\pm i\sqrt{1-\frac{k^2}{4}}$, вещественные части отрицательны при $k>0$. Следовательно, нулевое решение асимптотически устойчиво, что соответствует затухающим колебаниям.
\end{example}

\end{document}
